\scanpage{22}
Extensions quadratiques.

1. Un \deleted{\unsure{schéma}}\inserted{revêtement} quadratique d'un schéma $S$ est un
$S$-schéma $\Omega$ fini loc. libre de rang 2 sur $S$. La catégorie
de ces revêtements quadratiques (sur $S$ fixé) est
équivalente à celle des $\underline{O}_S$-Algèbres (comm. unifères)
\inserted{$B$} loc. libres de rang 2, ou \inserted{dites aussi} extensions quadratiques
de $\underline{O}_S$. Pour $S$ variable, on trouve une catégorie
fibrée sur $(\mathrm{Sch})_{/S}$
\[
\Omega=\operatorname{Spec}(B) \tag*{(1)}
\]
Soit \deleted{\unsure{le}}
\[
\check L=B/\underline{O}_S\qquad 0\mathbin{-}\underline{O}_S\xrightarrow{i}B\mathbin{-}\check L\longrightarrow0 \tag*{(2)}
\]
module inversible. On a une structure d'extension
\[
0\mathbin{-}L^{\deleted{\unsure{\vee}}}\longrightarrow\check B\xrightarrow{\pi}\underline{O}_S\mathbin{-}0 \tag*{(3)}
\]
i.e. un fibré affine de \unsure{rg|rang} 1 sur $S$, ayant
$L^{\deleted{\unsure{\vee}}}$ comme module des translations, \deleted{les splittings}\deleted{scindages}
soit
\[
\begin{aligned}
 E&\subset\check{\mathbb V}(\check B)\simeq\mathbb V(B)=\operatorname{Spec}(\underline{\mathrm{Sym}}_{\underline{O}_S}(B))\\
 \Vert&\\[-3pt]
 \operatorname{Spec}(&\underline{\mathrm{Sym}}_{\underline{O}_S} B/(e_B-1))
\end{aligned} \tag*{(4)}
\]
où $e$ désigne la section unité de $B$. Les sections de
$E$ sur $S$ correspondent aux scindages de (3), ou
aux hom de $\underline{O}_S$-modules $B\longrightarrow\underline{O}_S$ qui
transforment $e_B$ en $1_S$, i.e. inverses à gauche de $i$.
Parmi ceux-ci il y a les \deleted{scindages} hom de
$\underline{O}_S$-Algèbres. Donc (\unsure{notant|aidant} que l'application inj.
\[
\deleted{\unsure{f(S)}}\ \Omega(S)\hookrightarrow E(S) \tag*{(5)}
\]
ainsi définie commute à tt chgt de base $S'\longrightarrow S)$

% archive sheet number: 21

